% =====================================================================
%  sections/06_implementation.tex  --  VI. IMPLEMENTATION
% =====================================================================
\clearpage
\section{IMPLEMENTATION}
\label{sec:implementation}

% ---------------------------------------------------------------------
\subsection{Application Flow}
\label{sub:application-flow}

This diagram traces a single user query end to end through the running
system, showing the concrete sequence of steps behind the AI
architecture described earlier. It follows the request from language
detection and vector store retrieval through LLM query execution to
the final explanation returned to the user.

\projfigure{0.95}{fig-application-flow.png}{\protect\projnameplain{}'s end-to-end query processing flow, from user query and language detection through vector store retrieval, LLM query execution, and final result explanation.}{fig:application-flow}

% ---------------------------------------------------------------------
\subsection{Implementation Challenges and Solutions}
\label{sub:implementation-challenges-and-solutions}

Three challenges shaped the way the system was built.

\subsubsection*{1. The database schema was too large for a single prompt}
Our first approach placed the whole schema in the prompt so the model
could see every table and column. The prompt quickly grew past a
practical size, responses slowed down, and accuracy dropped because the
relevant tables were buried among unrelated ones.

\noindent \textbf{How we solved it.} We embedded the table and column descriptions
into a Chroma vector store and retrieve only the fragments that match
the current question. The prompt now carries a small, relevant slice of
the schema instead of all of it, which shortened response time and
improved the quality of the generated SQL.

\subsubsection*{2. Questions arrive in two languages}
The Kyungrinara dataset uses Korean labels while most of our team and
our schema documentation are in English. Questions asked in Korean did
not match the stored schema descriptions, so retrieval returned the
wrong tables.

\noindent \textbf{How we solved it.} A detection step identifies Korean by
checking the message against the Hangul character range, falling back to
a language-detection library only for other scripts. Korean questions
are translated to English before retrieval, so one English schema index
serves both languages instead of maintaining two.

\subsubsection*{3. Generated SQL is not always valid}
The model sometimes returned a statement wrapped in markdown fences,
written in a dialect PostgreSQL does not accept, or truncated with
unbalanced parentheses. A query could also be perfectly valid yet return
no rows, which left the interface with an empty chart and no
explanation.

\noindent \textbf{How we solved it.} Every generated statement passes through a
cleaning and validation step that strips markdown fences, converts
dialect differences, and rejects a query missing a \code{SELECT} clause
or with unbalanced parentheses. Failures are fed back to the model as
error context for a bounded number of retries, and an empty result is
tagged so the same handler reports it to the user in words rather than
showing a blank chart.